\section{Product Motivation}

\begin{frame}{Predictions Are Not Decisions}
  \centering
  \begin{tabular}{p{0.43\textwidth}p{0.43\textwidth}}
    \toprule
    \textbf{Supervised learning predicts} & \textbf{RL chooses} \\
    \midrule
    Truck arrival time & Which truck to unload next \\
    Unloading duration & Which pit to assign \\
    Grain quality & Which silo to route into \\
    Train delay & How to schedule train loading \\
    \bottomrule
  \end{tabular}

  \vspace{0.5cm}
  \takeaway{Supervised learning predicts what happens next. Reinforcement learning chooses what should happen next.}
  \note{Open with the grain terminal. Prediction is useful, but the product tension starts when the system must choose an action that changes the operation.}
\end{frame}

\begin{frame}{Actions Change the Future}
  \centering
  \begin{tikzpicture}[node distance=0.55cm,>=Latex,scale=0.94,transform shape]
    \node[draw,rounded corners,fill=Panel,minimum width=2.0cm,minimum height=0.8cm] (queue) {truck queue};
    \node[draw,rounded corners,fill=Panel,right=of queue,minimum width=2.0cm,minimum height=0.8cm] (pit) {unloading pit};
    \node[draw,rounded corners,fill=Panel,right=of pit,minimum width=2.1cm,minimum height=0.8cm] (route) {conveyor route};
    \node[draw,rounded corners,fill=Panel,right=of route,minimum width=2.0cm,minimum height=0.8cm] (silo) {silo / blend};
    \node[draw,rounded corners,fill=Panel,right=of silo,minimum width=2.0cm,minimum height=0.8cm] (train) {train loading};
    \draw[->,thick] (queue) -- (pit);
    \draw[->,thick] (pit) -- (route);
    \draw[->,thick] (route) -- (silo);
    \draw[->,thick] (silo) -- (train);
    \draw[->,thick,SignalGreen] (train.north) .. controls +(0,0.95) and +(0,0.95) .. node[above]{new state} (queue.north);
  \end{tikzpicture}

  \vspace{0.45cm}
  Choosing a truck, pit, route, or silo changes future queues, storage, blending options, and train schedules.

  \vspace{0.25cm}
  \productbox{In RL, the model does not just observe the data distribution. It helps create the next data distribution.}
  \note{This is the main conceptual gap. Supervised learning assumes a fixed data distribution; a policy changes the states it will see next.}
\end{frame}

\begin{frame}{States, Actions, and Policies}
  \centering
  \begin{tabular}{p{0.22\textwidth}p{0.65\textwidth}}
    \toprule
    \textbf{Object} & \textbf{Grain-terminal example} \\
    \midrule
    state $s_t$ & current queue, silo inventory, grain quality, train progress \\
    action $a_t$ & unload truck, choose pit, route conveyor, assign silo, fill car \\
    policy $\pi(a\mid s)$ & rule or model that chooses actions from states \\
    next state $s_{t+1}$ & the facility after the action has changed it \\
    \bottomrule
  \end{tabular}

  \vspace{0.55cm}
  \takeaway{A decision problem starts with states, actions, and a policy that maps situations to decisions.}
  \note{Define the basic vocabulary before imitation learning uses it. State means the information available at decision time. Action means the decision we control. A policy is the rule or learned model that maps states to actions.}
\end{frame}

\begin{frame}{Why RL Needs More Than Logs}
  \centering
  \begin{tikzpicture}[node distance=0.9cm,>=Latex]
    \node[draw,rounded corners,fill=Panel,minimum width=3.3cm,minimum height=0.9cm] (sim) {simulator available};
    \node[draw,rounded corners,fill=Panel,right=2.0cm of sim,minimum width=3.3cm,minimum height=0.9cm] (logs) {logs only};
    \node[below=0.35cm of sim,text width=3.5cm,align=center] {try novel actions\\estimate consequences};
    \node[below=0.35cm of logs,text width=3.5cm,align=center] {improve cautiously\\stay near support};
    \draw[->,thick,SignalGreen] (sim.south) -- ++(0,-0.6);
    \draw[->,thick,SignalGreen] (logs.south) -- ++(0,-0.6);
  \end{tikzpicture}

  \vspace{1.0cm}
  \takeaway{RL is about policy improvement, not just action prediction.}
  \note{This sets up the later offline RL material. If we have a simulator, RL can evaluate counterfactual actions. If we only have logs, offline RL must improve while respecting dataset support.}
\end{frame}

\begin{frame}{The RL Question}
  \Large
  Given logged grain-terminal decisions and delayed outcomes, can we learn a better policy without risky online experimentation?

  \vspace{0.6cm}
  \normalsize
  Logged data:
  \[
    (s_t, a_t, \text{outcome}, s_{t+1})
  \]

  \vspace{0.2cm}
  \productbox{This is the bridge from product operations to imitation learning and the RL formalism.}
  \note{Use this as the transition into imitation learning and then the RL formalism. We now have states, actions, delayed outcomes, and a policy to improve.}
\end{frame}
